\subsection*{Relative attention and the local-window limit}
For a centered position $i$, the original question encoder uses the neighborhood $q_{i-4:i+4}$. Its two convolutions have dilation one and three. After encoding, question pooling and question-to-passage attention are symmetric under a permutation of the resulting features. For a given output prefix, a question position's copy bias depends on at most three preceding source cells. Those cells are contained in the same neighborhood. With the same boundary context, passage and multiset of nine-cell windows, each distinct emitted image therefore receives the same total probability in exact arithmetic. Induction extends the result to a deterministic complete output. Four such requests with distinct target streams permit at most one exact response per group from this architecture.

We compared two residual global attention layers with four heads. Both update a question embedding by a weighted sum of projected question values. In the ordered layer head $h$ adds $-|j-i|/\tau_h-8\,\mathbf{1}\{d_h(j-i)<0\}$ to the scaled dot-product score from position $i$ to $j$, with directions $d_h$ in $(-1,-1,+1,+1)$ and length scales $\tau_h$ in $(12,48,12,48)$. The unordered layer uses zero positional bias. Each residual output projection starts at zero. Both fits start from the English-matched writer and receive the same images, answer likelihood, replay, initialization and optimizer schedule. We vary which clause comes first and which field it mentions first. The two test frames each contribute four layouts for 284 headwords, with the complete nine-cell multiset checked on every four-intent group. The historical target appears only in the source evidence and is unavailable in the fixed writing memory.
